Table~7 gives the characters of the representations of these groups. Isomorphic groups possess the same representations. They are therefore placed together in a single table. In the first rows, the numerals preceding the symbols for group elements give the numbers of elements in the corresponding classes (see \S~93). The first columns contain the adopted conventional labels for the representations. The letters $A$ and $B$ label one-dimensional representations. The letter $E$ labels two-dimensional representations. The letter $F$ labels three-dimensional representations. (The symbol $E$ for a two-dimensional irreducible representation must not be confused with the symbol $E$ for the identity element of the group!)\footnote{The reason why two complex-conjugate one-dimensional representations are denoted as a single two-dimensional representation will become clear in \S~96.} The basis functions of representations $A$ are symmetric under rotations about the principal axis of order $n$. The basis functions of representations $B$ are antisymmetric under these rotations. For functions having different symmetry under reflection $\sigma_h$, the labels differ in the number of primes, one or two. The subscripts $g$ and $u$ indicate symmetry with respect to inversion. Alongside the representation labels, the letters $x,y,z$ indicate the representations according to which the coordinates themselves transform. The $z$ axis is always chosen along the principal symmetry axis. The letters $\varepsilon$ and $\omega$ denote:
